% Scheletro di lavoro, non testo definitivo.
% Le coordinate indicate sono piste di lettura da verificare sul testo greco prima della stesura.
% Metodo: testo originale, traduzione, ricostruzione argomentativa, analisi critica.
\chapter{La giustizia nella Repubblica}
\label{cap:platone:giustizia}

% Problema: Perché le prime definizioni risultano insufficienti?
% Fonti da verificare: Repubblica I 327a–354c.
\section{La domanda sulla giustizia}

\subsection{Cefalo e Polemarco: debito, verità e amicizia}

\subsection{Trasimaco: la giustizia come utile del più forte}

% Problema: Perché essere giusti quando l'ingiustizia sembra conveniente?
% Fonti da verificare: Repubblica II 357a–367e.
\section{La sfida di Glaucone e Adimanto}

\subsection{L'anello di Gige e la giustizia come convenzione}

\subsection{Il problema della giustizia per se stessa}

% Problema: Perché il problema dell'anima viene affrontato attraverso la città?
% Fonti da verificare: Repubblica II 368c–376e; IV 427d–434c.
\section{La città come ipotesi di ricerca}

\subsection{Dall'individuo alla città}

\subsection{Virtù della città e definizione della giustizia}

% Problema: In che senso ciascuna parte svolge il proprio compito?
% Fonti da verificare: Repubblica IV 435b–445e; raccordo con cap. 27.
\section{La giustizia nell'anima}

\subsection{La corrispondenza fra città e anima}

\subsection{Giustizia e ordine interiore}
